% ===========================================================================
% hebrew_pages.tex -- the Hebrew part of the thesis (XeLaTeX + biber).
%
% It is a SEPARATE document. It reads the English build's labels through
% xr-hyper (prefix en-) so that its table of contents shows the English
% chapter/section numbers and page numbers, continues the English page
% numbering, and links into the English PDF by name; merge_thesis.py turns
% those names into working links after the two PDFs are concatenated.
%
% Fonts: Linux Libertine O (Latin), David CLM (Hebrew), Latin Modern Roman
% (Latin text inside Hebrew paragraphs). All three ship with TeX Live; on
% macOS run ./compile_thesis.sh install-fonts once. See README.md.
% ===========================================================================
\documentclass{article}

\usepackage{xcolor}            % must be loaded before bidi (which polyglossia pulls in)
\usepackage{xr-hyper}          % read labels from build/main.aux (prefix en-)
\usepackage{amsmath}
\usepackage[urlcolor=blue!50!black,linkcolor=blue!50!black,citecolor=green!40!black,colorlinks]{hyperref}
\usepackage{graphicx}
\usepackage{polyglossia}
\usepackage{fontspec}
\usepackage{amssymb}
\usepackage{setspace}
\setstretch{1.25}
\usepackage{etoolbox}

% Same page geometry as main.tex (bottom margin differs slightly: the Hebrew
% pages have no running header, so the text block can be a little taller).
\usepackage[a4paper,twoside,inner=24mm,outer=24mm,top=28mm,bottom=22mm]{geometry}

% The SAME bibliography settings and .bib file as main.tex. Alphabetic labels
% are derived from author and year, so [Kar72] here is [Kar72] in the English
% bibliography, and the citation links resolve to the English entries after
% the merge.
\usepackage[backend=biber,natbib=true,style=alphabetic,sorting=none,
            maxcitenames=2,maxbibnames=50,doi=true,url=true,isbn=false,eprint=false]{biblatex}
\renewcommand*{\multicitedelim}{\addcomma\space}
\addbibresource{bibliography/references.bib}

\setmainfont{Linux Libertine O}
\setdefaultlanguage{hebrew}    % headings and paragraphs are right-to-left
\setotherlanguage{english}
\newfontfamily\hebrewfont[Script=Hebrew,Scale=MatchLowercase]{David CLM}
% polyglossia needs a Hebrew monospace font as soon as \texttt appears in Hebrew text
\newfontfamily\hebrewfonttt[Script=Hebrew,Scale=MatchLowercase]{David CLM}

\usepackage{bidi}

% Latin text inside a Hebrew paragraph: keep it left-to-right and in a Latin
% font.  \LR comes from bidi.  \en{...} is for words, \hcite{...} for citations.
\newcommand{\en}[1]{\beginL{\fontspec{Latin Modern Roman}#1}\endL}
\newcommand{\hcite}[1]{\mbox{\LR{{\fontspec{Latin Modern Roman}\cite{#1}}}}}

% ---- Cross-document references into the English build ----
% The build script compiles English first, so build/main.aux exists.
\IfFileExists{build/main.aux}{\externaldocument[en-]{build/main}}{}
% Numbers and pages of English headings, printed left-to-right and WITHOUT
% their own link (the surrounding \hyperlink provides the link; merge_thesis.py
% needs exactly one destination per link).
\newcommand{\englishTOCpage}[1]{\LR{\pageref*{en-#1}}}
\newcommand{\englishTOCnumber}[1]{\LR{\ref*{en-#1}}}
% A clickable reference to an English theorem/figure/definition that prints its number:
\newcommand{\enref}[2]{\hyperlink{target:#1}{\englishTOCnumber{#2}}}

% ---------------------------------------------------------------------------
% thesis_info.tex -- the only file you must edit to make the thesis yours.
% It is read by BOTH documents (main.tex and hebrew_pages.tex), so the cover
% pages, the Hebrew cover and the abstracts all agree.
% ---------------------------------------------------------------------------

% Title
\newcommand{\thesistitle}{A Thesis Template for Ben-Gurion University}
\newcommand{\thesistitlehe}{תבנית לעבודת גמר באוניברסיטת בן-גוריון}

% Author
\newcommand{\thesisauthorname}{Israel Israeli}
\newcommand{\thesisauthornamehe}{ישראל ישראלי}

% Supervisor(s). For two supervisors, define \thesissecondsupervisername and
% \thesissecondsupervisernamehe as well and adjust Main_pages/cover_page.tex,
% Main_pages/title_page.tex and the Hebrew cover at the end of hebrew_pages.tex.
\newcommand{\thesissupervisername}{Prof. Max Musterman}
\newcommand{\thesissupervisernamehe}{פרופ' מקס מוסטרמן}

% Faculty and department
\newcommand{\thesisfaculty}{The Faculty of Engineering}
\newcommand{\thesisfacultyhe}{הפקולטה להנדסה}
\newcommand{\thesisdepartment}{The Department of Industrial Engineering and Management}
\newcommand{\thesisdepartmenthe}{המחלקה להנדסת תעשייה וניהול}

% Degree statement (PhD wording; change for an MSc thesis)
\newcommand{\thesisdegreetext}{Thesis submitted in partial fulfillment of the requirements\\for the degree of Doctor of Philosophy}
\newcommand{\thesisdegreetexthe}{חיבור זה מהווה חלק מהדרישות לקבלת התואר דוקטור לפילוסופיה}

% Submission date
\newcommand{\thesismonth}{September}
\newcommand{\thesismonthhe}{ספטמבר}
\newcommand{\thesisyear}{2026}


\begin{document}
% Continue the Arabic page numbering of the English part (LastPage is
% defined there by the lastpage package).
\setcounter{page}{\numexpr\getpagerefnumber{en-LastPage}+1\relax}
\pagestyle{empty}
\selectlanguage{hebrew}

% ---------------------------------------------------------------------------
% Hebrew table of contents. One line per English entry:
%   \englishTOCnumber{<label>}  the chapter/section number from the English build
%   \hyperlink{<anchor>}{...}    the clickable Hebrew title
%   \englishTOCpage{<label>}    the page number from the English build
% Chapters: anchor = label = chapter:<key>. Sections: anchor target:sec:<key>,
% label toc:sec:<key>. Unnumbered chapters have no number.
% ---------------------------------------------------------------------------
\begin{center}
{\Huge\bfseries תוכן עניינים}
\end{center}
\vspace{1.2em}
\noindent
\begin{flushright}
{\Large
\englishTOCnumber{chapter:intro}.\ \hyperlink{chapter:intro}{מבוא}
\dotfill \englishTOCpage{chapter:intro}\\[0.2em]
\hspace{2em}\englishTOCnumber{toc:sec:what}\ \hyperlink{target:sec:what}{מה התבנית נותנת}
\dotfill \englishTOCpage{toc:sec:what}\\[0.1em]
\hspace{2em}\englishTOCnumber{toc:sec:quickstart}\ \hyperlink{target:sec:quickstart}{התחלה מהירה}
\dotfill \englishTOCpage{toc:sec:quickstart}\\[0.2em]
\englishTOCnumber{chapter:build}.\ \hyperlink{chapter:build}{תהליך הבנייה}
\dotfill \englishTOCpage{chapter:build}\\[0.2em]
\hspace{2em}\englishTOCnumber{toc:sec:script}\ \hyperlink{target:sec:script}{מה הסקריפט עושה}
\dotfill \englishTOCpage{toc:sec:script}\\[0.1em]
\hspace{2em}\englishTOCnumber{toc:sec:settle}\ \hyperlink{target:sec:settle}{מדוע שני השלבים הראשונים חוזרים}
\dotfill \englishTOCpage{toc:sec:settle}\\[0.1em]
\hspace{2em}\englishTOCnumber{toc:sec:troubleshooting}\ \hyperlink{target:sec:troubleshooting}{גופנים, אוברליף ופתרון תקלות}
\dotfill \englishTOCpage{toc:sec:troubleshooting}\\[0.2em]
\englishTOCnumber{chapter:links}.\ \hyperlink{chapter:links}{עוגנים וקישורים בין שני החלקים}
\dotfill \englishTOCpage{chapter:links}\\[0.2em]
\hspace{2em}\englishTOCnumber{toc:sec:anchors}\ \hyperlink{target:sec:anchors}{עוגנים ותוויות}
\dotfill \englishTOCpage{toc:sec:anchors}\\[0.1em]
\hspace{2em}\englishTOCnumber{toc:sec:sections}\ \hyperlink{target:sec:sections}{פרקים וסעיפים}
\dotfill \englishTOCpage{toc:sec:sections}\\[0.1em]
\hspace{2em}\englishTOCnumber{toc:sec:objects}\ \hyperlink{target:sec:objects}{משפטים, איורים, טבלאות והגדרות}
\dotfill \englishTOCpage{toc:sec:objects}\\[0.1em]
\hspace{2em}\englishTOCnumber{toc:sec:merge}\ \hyperlink{target:sec:merge}{שלב המיזוג}
\dotfill \englishTOCpage{toc:sec:merge}\\[0.2em]
\englishTOCnumber{chapter:hebrew}.\ \hyperlink{chapter:hebrew}{כתיבת העמודים העבריים}
\dotfill \englishTOCpage{chapter:hebrew}\\[0.2em]
\hspace{2em}\englishTOCnumber{toc:sec:rtl}\ \hyperlink{target:sec:rtl}{טקסט מימין לשמאל וקטעים לטיניים}
\dotfill \englishTOCpage{toc:sec:rtl}\\[0.1em]
\hspace{2em}\englishTOCnumber{toc:sec:hcite}\ \hyperlink{target:sec:hcite}{ציטוטים}
\dotfill \englishTOCpage{toc:sec:hcite}\\[0.1em]
\hspace{2em}\englishTOCnumber{toc:sec:enref}\ \hyperlink{target:sec:enref}{הפניות לתוצאות באנגלית}
\dotfill \englishTOCpage{toc:sec:enref}\\[0.1em]
\hspace{2em}\englishTOCnumber{toc:sec:hebrewtoc}\ \hyperlink{target:sec:hebrewtoc}{תוכן העניינים העברי}
\dotfill \englishTOCpage{toc:sec:hebrewtoc}\\[0.2em]
\hyperlink{chapter:conclusion}{מסקנות}
\dotfill \englishTOCpage{chapter:conclusion}\\[0.2em]
\hyperlink{chapter:bib}{ביבליוגרפיה}
\dotfill \englishTOCpage{toc:bibliography}\\[0.2em]
\hyperlink{chapter:hebrew:summary}{תקציר}
\dotfill \LR{\pageref*{toc:hebrew:summary}}\\[0.2em]
\hyperlink{chapter:hebrew:abstract}{תמצית}
\dotfill \LR{\pageref*{toc:hebrew:abstract}}\\[0.2em]
}
\end{flushright}
\newpage

% ---------------------------------------------------------------------------
% Hebrew summary (תקציר). The \hypertarget is what the English TOC and the
% Hebrew TOC link to; the \label is where both TOCs read its page number.
% ---------------------------------------------------------------------------
\noindent
{\huge\textbf{תקציר}}
\vspace{.5cm}
\hypertarget{chapter:hebrew:summary}{}\label{toc:hebrew:summary}%

\large
חיבור זה הוא תבנית לעבודת גמר באוניברסיטת בן-גוריון בנגב, והוא גם המדריך לשימוש בה.
החלק האנגלי נכתב במסמך אחד ומקומפל עם
\en{pdfLaTeX},
והעמודים העבריים, תוכן העניינים, התקציר, התמצית והכריכה העברית, נכתבים במסמך שני ומקומפלים עם
\en{XeLaTeX}.
הסקריפט
\en{compile\_thesis.sh}
בונה את שניהם, חוזר על הבנייה עד ששני המסמכים קוראים זה מזה את מספרי העמודים הנכונים, ומאחד אותם לקובץ אחד
(ראו איור
\enref{fig:pipeline}{fig:pipeline}).

\textbf{ציטוטים.}
ציטוט בתוך טקסט עברי נכתב בפקודה
\en{\texttt{\textbackslash hcite}}
בדיוק כמו
\en{\texttt{\textbackslash cite}}:
לדוגמה, בעיות רבות הן
\en{NP}-קשות
\hcite{Karp72,GJ79},
ותורת הסיבוכיות הפרמטרית
\hcite{downey2013fundamentals,cygan2015parameterized}
מציעה דרך להתמודד איתן.
התוויות בסוגריים זהות לאלו שברשימה הביבליוגרפית האנגלית, והלחיצה עליהן מובילה לערך המתאים שם.

\textbf{עוגנים.}
כל אובייקט בחלק האנגלי שנושא עוגן בשם
\en{\texttt{target:...}}
ניתן לקישור מכאן (הגדרה
\enref{def:anchor}{def:anchor}).
הפקודה
\en{\texttt{\textbackslash enref}}
מדפיסה את מספרו של האובייקט, כפי שנקרא מהבנייה האנגלית, והופכת אותו לקישור:
כך, המשפט המרכזי של התבנית הוא משפט
\enref{thm:main}{thm:main},
ומוסכמות השמות מסוכמות בטבלה
\enref{tab:naming}{tab:naming}.
המספרים בשורות אלו אינם כתובים בקובץ העברי, ולכן מספור מחדש של החלק האנגלי אינו שובר אותם.

\textbf{תוכן העניינים.}
כל שורה בתוכן העניינים העברי בנויה משלושה רכיבים: מספר הפרק או הסעיף, הכותרת העברית כקישור, ומספר העמוד.
המספר והעמוד נקראים מהבנייה האנגלית; הכותרת והסדר הם באחריות הכותב.
\newpage

% ---------------------------------------------------------------------------
% Hebrew abstract (תמצית): the Hebrew counterpart of Main_pages/abstract.tex.
% ---------------------------------------------------------------------------
\setstretch{1.45}
\noindent
{\huge\textbf{תמצית}}
\vspace{.5cm}
\hypertarget{chapter:hebrew:abstract}{}\label{toc:hebrew:abstract}%

\large
מסמך זה הוא תבנית לעבודת גמר באוניברסיטת בן-גוריון בנגב, שבה חלק אנגלי וחלק עברי מופקים כקובץ אחד שהקישורים בו פועלים בשני הכיוונים.
החלק האנגלי מסביר את תהליך הבנייה, את מוסכמות העוגנים והתוויות המקשרות בין שני החלקים, ואת הפקודות לכתיבת העמודים העבריים, ומדגים בדרך כל רכיב שעבודה אמיתית זקוקה לו: פרקים וסעיפים, משפט המנוסח במבוא ומוכח בהמשך, הגדרה, איור, טבלה, אלגוריתם, קטעי קוד וציטוטים.
החליפו תמצית זו בתמצית של העבודה שלכם.
\newpage

% ---------------------------------------------------------------------------
% Hebrew cover page: the LAST page of the bound thesis (read from the back).
% Its fields come from thesis_info.tex.
% ---------------------------------------------------------------------------
\begin{titlepage}
    \begin{center}
        \vspace*{1cm}
        \includegraphics[width=0.1\textwidth]{Main_pages/bgu.png}\\
        \textbf{אוניברסיטת בן-גוריון בנגב}\\
        \thesisfacultyhe\\
        \thesisdepartmenthe

        \vspace{2cm}

        {\Huge \textbf{\thesistitlehe}}

        \vspace{1.5cm}

        \thesisdegreetexthe~\en{(PhD)}

        \vspace{1cm}

        \textbf{\thesisauthornamehe}

        \vspace{1cm}

        בהנחיית \textbf{\thesissupervisernamehe}

        \vfill

        \thesismonthhe \quad \en{\thesisyear}
    \end{center}
\end{titlepage}

\end{document}
